\documentclass[11pt]{article}

\usepackage[a4paper,left=28mm,right=28mm,top=25mm,bottom=25mm]{geometry}
\usepackage[T1]{fontenc}
\usepackage{lmodern}
\usepackage{amsmath,amssymb,amsthm,mathtools}
\usepackage[hidelinks]{hyperref}

\newtheorem{theorem}{Theorem}
\newtheorem{lemma}[theorem]{Lemma}

\newcommand{\E}{\mathbb E}
\newcommand{\Pp}{\mathbb P}
\newcommand{\R}{\mathbb R}
\newcommand{\dd}{\,\mathrm d}
\newcommand{\abs}[1]{\lvert#1\rvert}
\newcommand{\pos}[1]{[#1]_{+}}

\title{A \(0.4395\) Upper Bound for the i.i.d. Berry--Esseen Constant}
\author{%
  Haonan Xiao\({}^{1}\) \qquad Chunlin Li\({}^{2}\)\\[3pt]
  {\small \({}^{1}\)Department of Statistics, Iowa State University}\\[-1pt]
  {\small \({}^{2}\)Department of Statistics, University of Virginia}
}
\date{September 2026}

\hypersetup{
  pdftitle={A 0.4395 Upper Bound for the i.i.d. Berry-Esseen Constant},
  pdfauthor={Haonan Xiao and Chunlin Li},
  pdfsubject={A formally verified upper bound in the classical i.i.d. Berry-Esseen inequality}
}

\begin{document}
\maketitle

\begingroup
\renewcommand{\thefootnote}{}
\footnotetext[1]{Emails:
  \href{mailto:xhn@iastate.edu}{\texttt{xhn@iastate.edu}} (H.~Xiao) and
  \href{mailto:chunlin@virginia.edu}{\texttt{chunlin@virginia.edu}} (C.~Li).}
\endgroup

\begin{abstract}
We prove that the absolute constant in the classical Berry--Esseen inequality
for identically distributed summands is at most \(879/2000=0.4395\), uniformly
over all sample sizes and all centered variance-one laws with finite absolute
third moment. The proof couples the symmetrized third-moment ratio
\(r=\E|X-X'|^3/(2\E|X|^3)\) with \(d=1-\E|X|\), where \(X'\) is an
independent copy of \(X\). Their joint constraints give a characteristic-function
comparison with the symmetric two-point law and a stronger modulus estimate
when \(d\) is small. Prawitz smoothing reduces the remaining argument to
continuous-parameter inequalities, proved by exact outward-rounded dyadic
certificates. The analytic reduction and the exact certificates are formally
verified in Lean. The proof was developed autonomously using OpenAI Codex,
primarily with GPT-5.6 Sol.
\end{abstract}
\section{Introduction}\label{sec:introduction}

Let \(X,X_1,X_2,\ldots\) be independent and identically distributed real
random variables with
\[
  \E X=0,\qquad \E X^2=1,\qquad
  \rho:=\E\abs{X}^3<\infty,
\]
and call a law satisfying these conditions admissible.  Set
\[
  S_n:=\frac{X_1+\cdots+X_n}{\sqrt n},\qquad
  \Delta_n:=\sup_{x\in\R}\abs{\Pp\{S_n\le x\}-\Phi(x)},
\]
where \(\Phi\) denotes the standard normal distribution function.  The
optimal i.i.d. Berry--Esseen constant is
\[
 C_0:=\inf\left\{C>0:\Delta_n\le
 C\frac{\rho}{\sqrt n}\text{ for all admissible laws and all }n\ge1\right\}.
\]
Esseen's classical example gives
\[
 C_0\ge C_{\mathrm E}:=
 \frac{\sqrt{10}+3}{6\sqrt{2\pi}}
 =0.409732\ldots\,\cite{esseen1956}.
\]
Shevtsova proved the published upper bound \(C_0\le0.4690\)
\cite{shevtsova2013}; Bobkov's survey records the resulting interval
\(0.4097<C_0<0.4690\) \cite[Theorem~20.1]{bobkov2016}.

\begin{theorem}\label{thm:main}
For every admissible law and every integer \(n\ge1\),
\[
  \Delta_n\le \frac{879}{2000}\,\frac{\rho}{\sqrt n}.
\]
Consequently, \(C_0\le0.4395\).
\end{theorem}

The starting point is the shared moment parameter
\[
 r:=\frac{\E|X-X'|^3}{2\rho},
\]
where \(X'\) is an independent copy of \(X\). Symmetrization yields a
modulus bound in terms of \(r\), while a stop-loss identity gives a bound
for the sine remainder \(\E[\sin(uX)-uX]\) in terms of the same \(r\).
Since \(\E X=0\), this remainder is the imaginary part of the
characteristic function. The resulting disk
comparison yields the \(0.45\) bound in the earlier version of this note.

The present improvement retains more of the moment information. If
\(\eta=\rho(r-1)\) and \(d=1-\E|X|\), then
\[
 0\le d\le\min(1-\rho^{-1},1-\eta).
\]
When these bounds force \(d\) to be small, expansion about \(|X|=1\)
controls the real and imaginary characteristic-function components
separately. This adds a comparison with \(\cos u\), the characteristic
function of the symmetric two-point law, without assuming that \(X\)
is symmetric. In the large-sample analysis it also gives an explicit
exponential modulus rate. Keeping the admissible range of the Gaussian
correction and the sample-size dependence in the telescoping bound further
sharpens the comparison. Write \(L=\rho/\sqrt n\).

A distribution-free bound handles \(L\ge1090/879\). Exact interval
certificates cover the remaining finite- and large-sample regimes,
including every boundary and arbitrarily small positive \(L\).

\section{Moment constraints for the characteristic function}\label{sec:moment-parameter}

The symmetrized third moment controls both the modulus and the imaginary
part of the characteristic function. Retaining the first absolute moment
gives an additional comparison with the symmetric two-point law.

\subsection{Stop-loss identities}\label{sec:moment-geometry}

The stop-loss representation of the symmetrized third moment restricts \(r\) and relates
it to the signed third moment, providing the real and imaginary constraints
used below.

Let \(X'\) be an independent copy of \(X\), and define
\[
 m:=\E\abs{X-X'}^3,\qquad r:=\frac{m}{2\rho}.
\]
For \(t\ge0\), introduce the stop-loss functions
\[
 A(t):=\E(X-t)_+,\qquad B(t):=\E(-X-t)_+,
\]
and put \(r_0:=\E X_+=\E X_-\).

\begin{lemma}\label{lem:stoploss}
One has
\begin{align}
 m&=2\rho+12\int_0^\infty\bigl(A(t)^2+B(t)^2\bigr)\dd t,
 \label{eq:m-stop}\tag{2.1}\\
 3\int_0^\infty A(t)^2\dd t&\le r_0\E X_+^2,
 \qquad
 3\int_0^\infty B(t)^2\dd t\le r_0\E X_-^2.
 \label{eq:hardy}\tag{2.2}
\end{align}
Consequently,
\[
 1\le r\le1+\frac{\E\abs X}{\rho}
 \le1+\frac1\rho\le2.
 \label{eq:r-range}\tag{2.3}
\]
If \(\mu:=\E X^3\), then
\[
 (r-1)^2+\left(\frac\mu\rho\right)^2\le1.
 \label{eq:moment-circle}\tag{2.4}
\]
\end{lemma}

\begin{proof}
For real \(x,y\),
\[
 \abs{x-y}^3=6\int_{\R}
 \bigl[(x-t)_+(t-y)_+ +(y-t)_+(t-x)_+\bigr]\dd t.
\]
Taking expectations, using independence, splitting the integral at zero,
and applying \(\E X=0\) gives
\[
 m=12\int_0^\infty\bigl(A^2+B^2+tA+tB\bigr)\dd t.
\]
Fubini's theorem gives
\(\int_0^\infty tA(t)\dd t=\E X_+^3/6\) and the analogous identity
for \(B\), proving \eqref{eq:m-stop}.

If \(0\le y\le z\), direct integration gives
\[
 3\int_0^\infty(y-t)_+(z-t)_+\dd t
 =\frac32y^2z-\frac12y^3
 \le\frac12(yz^2+y^2z).
\]
Averaging independent copies proves \eqref{eq:hardy}.  Since
\(\E X_+^2+\E X_-^2=1\), it follows that
\(m-2\rho\le4r_0=2\E\abs X\), and \eqref{eq:r-range} follows from
\(\E\abs X\le1\) and Lyapunov's inequality \(\rho\ge1\).

To prove \eqref{eq:moment-circle}, write
\[
 a_j:=\E X_+^j,\qquad b_j:=\E X_-^j\quad(j=2,3).
\]
Positive semidefiniteness of the covariance matrix of \((X_+,X_-)\),
together with \(X_+X_-=0\), gives \(a_2b_2\ge r_0^2\).
Cauchy's inequality gives \(a_2^2\le r_0a_3\) and
\(b_2^2\le r_0b_3\), hence \(r_0\le\sqrt{a_3b_3}\).  Therefore
\[
 m-2\rho\le4\sqrt{a_3b_3}
 =2\sqrt{\rho^2-\mu^2}.
\]
Division by \(2\rho\) proves \eqref{eq:moment-circle}.
\end{proof}

The range in \eqref{eq:r-range} will control the modulus of the characteristic
function. Inequality \eqref{eq:moment-circle} shows at the same time that a
larger value of \(r-1\) forces the normalized third moment to be smaller.  The
next result transfers this tradeoff to the sine remainder at every frequency.

\subsection{A bound for the sine remainder}

The same parameter \(r\) controls \(\E[\sin(uX)-uX]\) at every frequency.

\begin{lemma}[Bound for the sine remainder]\label{lem:sine-circle}
For every \(u\in\R\),
\[
 \abs{\E(\sin(uX)-uX)}
 \le\frac{\rho\abs u^3}{6}\sqrt{1-(r-1)^2}.
 \label{eq:sine-circle}\tag{2.5}
\]
\end{lemma}

\begin{proof}
It is enough to take \(u>0\).  Define
\[
 J:=\int_0^\infty t(A+B)\dd t=\frac\rho6,
 \quad E_0:=\int_0^\infty(A^2+B^2)\dd t,
 \quad H_u:=\frac1u\int_0^\infty\sin(ut)(A-B)\dd t.
\]
The point of this normalization is geometric.  For a centered two-point law,
the normalized quantities satisfy the unit-disk inequality
\(E_0^2+H_u^2\le J^2\).  Stop-loss functions are affine under mixtures, so
Minkowski's inequality carries this constraint from two-point laws to an
arbitrary centered law.  We now make that mixture explicit.

Let \(\alpha_+\) be the restriction
of the law of \(X\) to \((0,\infty)\), and let \(\alpha_-\) be the reflection
of its restriction to \(( -\infty,0)\).  Centering gives
\[
 \int a\,\alpha_+(\dd a)=\int b\,\alpha_-(\dd b)=r_0,
\]
and the variance-one assumption implies \(r_0>0\).
On \((0,\infty)^2\), define
\[
 \nu(\dd a,\dd b):=\frac{a+b}{r_0}\,
 \alpha_+(\dd a)\alpha_-(\dd b),
\]
and add mass \(\Pp\{X=0\}\) at a zero component.  For each \((a,b)\), let
\(X_{a,b}\) take the values \(a\) and \(-b\) with probabilities
\(b/(a+b)\) and \(a/(a+b)\).  A direct calculation shows that the law of
\(X\) is the \(\nu\)-mixture of these centered two-point laws and the zero
law.  In particular, their stop-loss functions mix linearly to \((A,B)\).

For one such component, let \(J_{a,b},E_{a,b},H_{a,b}\) denote the three
quantities above, and put \(x=ua\), \(y=ub\).  Direct calculation gives
\[
 \frac{E_{a,b}}{J_{a,b}}=\frac{2xy}{x^2+y^2},\qquad
 \frac{H_{a,b}}{J_{a,b}}
 =-\frac{6(y\sin x-x\sin y)}{xy(x^2+y^2)}.
\]
For \(z>0\),
\[
 \left|\frac{\dd}{\dd(z^2)}\frac{\sin z}{z}\right|
 =\frac{\abs{z\cos z-\sin z}}{2z^3}\le\frac16,
\]
because
\[
 \abs{z\cos z-\sin z}
 =\left|\int_0^z t\sin t\dd t\right|
 \le\int_0^z t^2\dd t=\frac{z^3}{3}.
\]
It follows that \(E_{a,b}^2+H_{a,b}^2\le J_{a,b}^2\). Applying Jensen's
inequality at each \(t\), together with linearity of the mixture, gives
\[
 E_0\le\int E_{a,b}\nu(\dd a,\dd b),\qquad
 H_u=\int H_{a,b}\nu(\dd a,\dd b),\qquad
 J=\int J_{a,b}\nu(\dd a,\dd b).
\]
Minkowski's inequality in \(\R^2\) now proves
\[
 \sqrt{E_0^2+H_u^2}
 \le\int\sqrt{E_{a,b}^2+H_{a,b}^2}\nu(\dd a,\dd b)
 \le J.
\]

Two integrations of the sine remainder yield
\(\E(uX-\sin(uX))=u^3H_u\).  By \eqref{eq:m-stop},
\(E_0/J=r-1\). Thus the displayed bound on \(E_0^2+H_u^2\) gives
\(\abs{H_u}\le(\rho/6)\sqrt{1-(r-1)^2}\), proving
\eqref{eq:sine-circle}.
\end{proof}

Thus \(r\) controls the sine remainder uniformly in the frequency.  Together
with the modulus estimate below, this gives a joint constraint on the complex
characteristic-function error.

\subsection{Constraints involving the first absolute moment}\label{sec:radial-defect}

Write
\[
 s:=\E|X|,\qquad d:=1-s,\qquad \eta:=\rho(r-1).
\]
The Hardy bounds \eqref{eq:hardy}, before replacing \(s\) by 1, give
\[
 m-2\rho\le4r_0(\E X_+^2+\E X_-^2)=2s,
 \qquad\text{hence}\qquad \eta\le s.
\]
Cauchy--Schwarz also gives
\[
 1=(\E|X|^2)^2\le(\E|X|)(\E|X|^3)=s\rho,\qquad s\le1.
\]
Consequently,
\[
 0\le d\le\bar d:=\min(1-\rho^{-1},1-\eta).
 \label{eq:defect}\tag{2.6}
\]
The identity \(\E(|X|-1)^2=2d\) explains the role of this additional
constraint: small \(d\) forces \(|X|\) to be close to 1 in mean square.
Expanding the characteristic function about \(|X|=1\) will therefore
give a comparison with the symmetric two-point law.

\section{Characteristic-function comparisons}\label{sec:cf-bounds}

Let \(f(u):=\E e^{iuX}\). The smoothing inequality requires bounds on
both \(|f(u)|\) and \(|f(u)-e^{-u^2/2}|\): the first controls the
decay of the characteristic function of the sum, and the second its
difference from the Gaussian characteristic function. Put
\[
 g(v):=\cos v-1+\frac{v^2}{2},\qquad
 \kappa:=\max_{v>0}\frac{g(v)}{v^3}.
\]
Exact Taylor enclosures and proved derivative-sign bounds give
\[
 0.09916191350<\kappa<0.09916191353
\]
and a unique maximizer \(\vartheta\) satisfying
\[
 3.99589567<\vartheta<3.99589570.
\]
We use Prawitz's piecewise characteristic-function minorant
\cite{prawitz1975cf}, at the symmetrized scale \(r\rho\):
\[
 q(v):=
 \begin{cases}
 \frac12-\kappa v,&0\le v\le\vartheta,\\[1mm]
 \dfrac{1-\cos v}{v^2},&\vartheta<v\le2\pi,\\[2mm]
 0,&v>2\pi.
 \end{cases}
 \label{eq:q-def}\tag{3.1}
\]

\begin{lemma}\label{lem:q}
The function \(q\) is convex and nonincreasing, and
\[
 q(v)\le\frac{1-\cos v}{v^2}\qquad(v>0).
\]
Consequently,
\[
 \abs{f(u)}^2\le1-2u^2q(r\rho\abs u).
 \label{eq:modulus}\tag{3.2}
\]
\end{lemma}

\begin{proof}
The definition of \(\kappa\) gives the minorant on the first interval.
At \(\vartheta\), equality and equality of the first derivatives hold.
For \(\psi(v)=(1-\cos v)/v^2\),
\[
 \psi''(v)=\frac{v^2\cos v-4v\sin v+6(1-\cos v)}{v^4}.
\]
The numerator is nonnegative on \([\vartheta,2\pi]\), and
\(\psi'(2\pi)=0\).  The exact sign checks are part of the certificate in
Section~\ref{sec:certificate}.  Hence the three pieces join with
nondecreasing slopes and \(q\) is convex and nonincreasing.

Let \(Y=X-X'\).  Since \(\E Y^2=2\), the measure \(Y^2\dd\Pp/2\) is a
probability measure.  Jensen's inequality gives
\begin{align*}
 1-\abs{f(u)}^2
 &=\E(1-\cos(uY))\\
 &\ge2u^2q\left(\abs u\frac{\E\abs Y^3}{\E Y^2}\right)
 =2u^2q(r\rho\abs u),
\end{align*}
which is \eqref{eq:modulus}.
\end{proof}

Estimate \eqref{eq:modulus} will control the powers of \(f\) in the
smoothing argument. We next combine \eqref{eq:sine-circle} with the
Taylor remainder bounds to obtain a one-step Gaussian comparison
by planar maximization, retaining the same parameter \(r\).

The sine-remainder bound gives an
imaginary half-width \(\beta(r)\), the ordinary third-order remainder gives a
disk of radius \(d_0\), and the Gaussian correction shifts the comparison
point along the real axis by \(\epsilon_\rho(c)\).  Set
\[
 \beta(r):=\frac16\sqrt{1-(r-1)^2},\qquad d_0:=\frac16,
\]
and, for \(c\ge0\), define
\[
 \epsilon_\rho(c):=
 \begin{cases}
 \displaystyle\frac{\rho^2}{c^3}
 \left(e^{-c^2/(2\rho^2)}-1+\frac{c^2}{2\rho^2}\right),&c>0,\\[2mm]
 0,&c=0.
 \end{cases}
\]
Let
\[
 \Xi(r):=\{0,\kappa\}\cup
 \left(\left\{\frac{r-1}{6}\right\}\cap[0,\kappa]\right)
\]
and
\[
 D(\rho,r,c)^2:=\max_{x\in\Xi(r)}
 \left\{(x-\epsilon_\rho(c))^2+
 \min\bigl(\beta(r)^2,d_0^2-x^2\bigr)\right\}.
 \label{eq:D-def}\tag{3.3}
\]
Thus \(D(\rho,r,c)\) is the greatest possible distance from the Gaussian
comparison point to the feasible Taylor-remainder region. The maximum can
be computed by checking the real parts in \(\Xi(r)\).

\begin{lemma}\label{lem:one-step}
For every \(u\in\R\), with \(c=\rho\abs u\),
\[
 \abs{f(u)-e^{-u^2/2}}\le\rho\abs u^3D(\rho,r,c).
 \label{eq:one-step}\tag{3.4}
\]
\end{lemma}

\begin{proof}
At \(u=0\) both sides are zero. For \(u\ne0\), set
\[
 z:=\frac{f(u)-(1-u^2/2)}{\rho\abs u^3},
\]
the definition of \(\kappa\), Lemma~\ref{lem:sine-circle}, and the
third-order Taylor remainder give
\[
 0\le\Re z\le\kappa,\qquad
 \abs{\Im z}\le\beta(r),\qquad \abs z\le d_0.
\]
Moreover,
\[
 \frac{e^{-u^2/2}-(1-u^2/2)}{\rho\abs u^3}
 =\epsilon_\rho(c).
\]
Thus the normalized left side of \eqref{eq:one-step} is the distance
from the nonnegative real point \(\epsilon_\rho(c)\) to a point in the
intersection of a vertical strip, a horizontal strip, and the disk
of radius \(d_0\).
For fixed real part \(x\), the largest squared imaginary part is
\(\min(\beta(r)^2,d_0^2-x^2)\).  The squared distance can change form
only at \(x=(r-1)/6\), and its maximum on each piece is attained at an
endpoint.  These are exactly the points in \(\Xi(r)\), proving
\eqref{eq:D-def} and \eqref{eq:one-step}.
\end{proof}

\subsection{Comparison with the symmetric two-point law}\label{sec:radial-cf}

The first absolute moment gives a second comparison, with \(\cos u\).
The following bounds express its error directly in terms of \(d\),
using the mean-square identity at the end of Section~\ref{sec:radial-defect}.

\begin{lemma}\label{lem:radial}
For every \(u\in\R\),
\[
\begin{aligned}
 |\Re f(u)-\cos u|&\le d\bigl(|u\sin u|+u^2\bigr),\\
 |\Im f(u)|&\le\sqrt{2d}\,|\sin u-u\cos u|+du^2.
\end{aligned}
\label{eq:radial-cf}\tag{3.5}
\]
\end{lemma}

\begin{proof}
Let \(Y=|X|\), and let \(A\in\{-1,1\}\) be the sign of \(X\), with \(A=1\)
when \(X=0\). Then \(X=AY\), \(\E AY=0\), and
\(\E(Y-1)^2=2d\). Hence
\[
 |\E A|=|\E[A(Y-1)]|\le\sqrt{2d}.
\]
Global Taylor remainder bounds give
\[
\begin{aligned}
 \cos(uY)&=\cos u-u\sin u(Y-1)+R_c,\\
 \sin(uY)&=\sin u+u\cos u(Y-1)+R_s,
 \qquad |R_c|,|R_s|\le\tfrac12u^2(Y-1)^2.
\end{aligned}
\]
Take expectations, using \(\E(Y-1)=-d\) and
\(\E[A(Y-1)]=-\E A\). This gives
\[
 \Re f(u)-\cos u=du\sin u+\E R_c,\qquad
 \Im f(u)=\E A(\sin u-u\cos u)+\E(AR_s).
\]
Both remainder expectations have absolute value at most \(du^2\).
\end{proof}

Intersect these real and imaginary bounds with the constraints in the proof
of Lemma~\ref{lem:one-step}. This bounds both
\(|f(u)|\) and \(|f(u)-\cos u|\). The latter gives the additional estimate
\[
 |f(u)^n-e^{-nu^2/2}|
 \le n|f(u)-\cos u|\max(|f(u)|,|\cos u|)^{n-1}
      +|\cos(u)^n-e^{-nu^2/2}|.
 \label{eq:rademacher}\tag{3.6}
\]
Thus comparison with the two-point law also bounds the error in the
normal approximation.

When \(r\) is close to 2, the moment constraints force \(d\) to be
small uniformly in \(\rho\). The resulting exponential modulus bound
will control the powers of \(f\) in the large-sample estimate.

\begin{lemma}\label{lem:radial-rate}
If \(r\ge19/10\) and \(0\le u\le5/6\), then
\[
 |f(u)|\le e^{-(39/100)u^2}.
 \label{eq:radial-rate}\tag{3.7}
\]
\end{lemma}

\begin{proof}
From \eqref{eq:defect},
\[
 d\le1-\max(\rho^{-1},\rho(r-1))\le1-\sqrt{r-1}.
\]
The last inequality follows by considering whether
\(\sqrt{r-1}\le\rho^{-1}\). Since
\((237/250)^2\le9/10\), we have \(d\le13/250\).
Put \(x=u^2\) and
\[
 P(x)=1-\frac{99}{250}x+
       \left(\frac1{24}-\frac{13}{1500}\right)x^2+
       \frac{13}{30000}x^3.
\]
Taylor bounds for sine and cosine applied to \eqref{eq:radial-cf} give
\[
 |\Re f(u)|\le P(x),\qquad
 |\Im f(u)|\le\frac{71}{500}x.
\]
For the first bound, \(1-u^2/2>0\) is a lower bound on \(\Re f(u)\).
For the second, use
\(0\le\sin u-u\cos u\le u^3/3\),
\(\sqrt{2d}\le81/250\), and \(u\le5/6\); indeed,
\((81/250)(5/6)/3+13/250=71/500\).

Let \(Q_3(x)=1-\frac{39}{50}x+
\frac12(\frac{39}{50}x)^2-\frac16(\frac{39}{50}x)^3\).
For \(0\le x\le25/36\),
\[
 Q_3(x)-P(x)^2-(71x/500)^2=x\,q_*(x),
\]
where
\[
\begin{aligned}
q_*(x)={}&\frac3{250}+\frac{3061}{50000}x
 -\frac{40367}{750000}x^2-\frac{3729}{5000000}x^3\\
 &-\frac{143}{5000000}x^4-\frac{169}{900000000}x^5.
\end{aligned}
\]
Here \(q_*''<0\), \(q_*(0)=3/250>0\), and
\[
 q_*(25/36)=\frac{98569207799}{3482851737600}>0.
\]
Concavity proves \(q_*\ge0\) throughout the interval. The alternating
exponential bound \(Q_3(x)\le e^{-(39/50)x}\), followed by a square root,
proves \eqref{eq:radial-rate}.
\end{proof}

\subsection{The range of the Gaussian correction}\label{sec:feasible-disk}

For \(r>1\), admissibility restricts \(1\le\rho\le1/(r-1)\).
At fixed \(c\ge0\), \(\epsilon_\rho(c)\) decreases with \(\rho\).
For \(c>0\), this follows by writing it as
\[
 \frac{e^{-y}-1+y}{2cy},\qquad y=\frac{c^2}{2\rho^2},
\]
since the derivative of \((e^{-y}-1+y)/y\) is
\([1-(1+y)e^{-y}]/y^2\ge0\). The case \(c=0\) is direct. Thus
\[
 \epsilon_{1/(r-1)}(c)\le\epsilon_\rho(c)\le\epsilon_1(c).
\]
The squared distance in \eqref{eq:D-def} is convex in the Gaussian correction.
Maximizing over \(\Xi(r)\) preserves convexity, so the two feasible
endpoints give
\[
 D(\rho,r,c)\le D_*(r,c):=
 \max\{D(1/(r-1),r,c),D(1,r,c)\}.
 \label{eq:Dstar}\tag{3.8}
\]
This is used only when \(r\ge19/10\). Otherwise we use
\(D_0(r,c)\), the bound obtained from the larger interval
\([0,\epsilon_1(c)]\) by evaluating the same squared distance at its endpoints.

\section{Smoothing and the \texorpdfstring{\(0.4395\)}{0.4395} bound}\label{sec:prawitz}

Set \(L:=\rho/\sqrt n\). When \(L\) is large, no characteristic-function
estimate is needed. The following bound uses only the first two moments.

\begin{lemma}\label{lem:universal}
If \(W\) is centered and has variance one, then
\[
 \sup_x|\Pp\{W\le x\}-\Phi(x)|\le\frac{109}{200}.
\]
\end{lemma}

\begin{proof}
At \(x=0\), the absolute error is at most \(1/2\). For \(a>0\),
Cantelli's inequality gives
\[
 \Pp\{W\ge a\}\le\frac1{1+a^2},\qquad
 \Pp\{W\le-a\}\le\frac1{1+a^2}.
\]
On each half-line one signed CDF error is at most \(1/2\). For the
other, use the displayed inequalities and, on the positive half-line,
\(\Pp\{W>a\}\le\Pp\{W\ge a\}\). By symmetry of \(\Phi\), it suffices
to bound \(1/(1+a^2)-\Phi(-a)\). For \(a\ge1\) this is at most \(1/2\).
For \(0<a\le1\), the bound \((2\pi)^{-1/2}<2/5\) gives
\[
 \frac1{1+a^2}-\Phi(-a)
 \le\frac1{1+a^2}-\frac12+\frac{2a}{5}\le\frac{109}{200}.
\]
The last inequality is equivalent to \(p(a):=9-80a+209a^2-80a^3\ge0\).
The following exact decompositions prove it on the two indicated intervals:
\[
\begin{aligned}
p(a)&=20a^2(1-4a)+\frac{(189a-40)^2+101}{189},
       &&0\le a\le\tfrac14,\\
p(a)&=\frac{13}{16}+(a-\tfrac14)
       \bigl[(a-\tfrac14)(169-80a)+\tfrac{19}{2}\bigr],
       &&\tfrac14\le a\le1.
\end{aligned}
\]
Every term on the right is nonnegative on its stated interval.
\end{proof}

Put \(C=879/2000\) and \(R=1090/879\). Since
\[
 CR=\frac{109}{200},\qquad
 \frac{56}{45}-R=\frac{58}{13185}>0,
\]
Lemma~\ref{lem:universal} proves Theorem~\ref{thm:main} when \(L\ge R\).
We certify the slightly larger complementary domain \(L\le56/45\)
using the characteristic-function bounds and Prawitz smoothing.

For \(0<t<1\), define the Prawitz kernel
\[
 \begin{aligned}
 K(t)&:=\frac{1-t}{2}
 +\frac{i}{2}\left((1-t)\cot(\pi t)+\frac1\pi\right),\\
 K(-t)&:=\overline{K(t)}.
 \end{aligned}
 \label{eq:kernel}\tag{4.1}
\]
The Prawitz smoothing inequality converts characteristic-function estimates
into a bound on Kolmogorov distance \cite{prawitz1972}.
The normalization below is equivalent to the filter normalization recorded in
\cite[Section~2]{pinelis2013}.

\begin{lemma}[Prawitz smoothing inequality]\label{lem:prawitz}
Let \(Y\) be an integrable real random variable, with distribution function
\(F\) and characteristic function \(h\).
For \(T>0\) and \(0<t_0<1\),
\begin{align}
 \sup_x\abs{F(x)-\Phi(x)}
 \le{}&2\int_0^{t_0}\abs{K(t)}
       \abs{h(Tt)-e^{-T^2t^2/2}}\dd t \notag\\
 &+2\int_{t_0}^1\abs{K(t)}\abs{h(Tt)}\dd t \notag\\
 &+2\int_0^{t_0}\left|K(t)-\frac{i}{2\pi t}\right|
       e^{-T^2t^2/2}\dd t \notag\\
 &+\frac1\pi\int_{t_0}^{\infty}e^{-T^2t^2/2}\frac{\dd t}{t}.
 \label{eq:prawitz}\tag{4.2}
\end{align}
\end{lemma}

\begin{proof}
Let \(J(x)=(\sin(\pi x)/(\pi x))^2\), with its removable value at zero,
and let \(H\) be the odd band-limited function in Prawitz's construction.
The pointwise estimate
\(\abs{\operatorname{sgn}(x)-H(x)}\le J(x)\) gives
\[
 B_-(x):=\frac{1-H(x)-J(x)}2
 \le \mathbf 1_{(-\infty,0]}(x)
 \le \frac{1-H(x)+J(x)}2=:B_+(x).
\]
The nonconstant Fourier transforms of \(B_+\) and \(B_-\) are supported on
\([-1,1]\); under the normalization in \eqref{eq:kernel}, they are \(K(t)\)
and \(-K(-t)\), respectively.

Apply these bounds to \(T(Y-x)/(2\pi)\), take expectations, and use Fourier
inversion.  The finite first moment removes the apparent singularity at zero:
\[
 \abs{h(u)-1}\le \abs u\,\E\abs Y,
\]
so the zero-frequency constant may be subtracted before passing to the
symmetric Fourier limit.  Subtracting the exact normal inversion formula,
splitting the frequencies at \(t_0\), and using conjugate symmetry gives the
first three terms in \eqref{eq:prawitz}; the remaining normal frequencies give
the last integral.  Applying the same argument to \(-Y\) supplies the opposite
CDF inequality and completes the absolute-value bound.
\end{proof}

We apply Lemma~\ref{lem:prawitz} to the law of \(S_n\), which is integrable
under the third-moment assumption.  With
\(L=\rho/\sqrt n\) as above, choose
\[
 T:=\frac{2\pi}{rL},\qquad t_0:=\frac{19}{50}.
 \label{eq:parameters}\tag{4.3}
\]
\subsection{The finite-sample envelope}

For \(t\ge0\), write \(u=2\pi t/(r\rho)\), \(B=e^{-u^2/2}\), and
\(N=N(t):=B^n=e^{-T^2t^2/2}\). Set
\(\kappa_+=9916191353/10^{11}\). Unlike the exact \(\kappa\) in
\(q\) and \(D\), \(\kappa_+\) is a certified rational upper bound,
used in the real strip below.

First intersect the component bounds from Section~\ref{sec:cf-bounds}.
With \(\bar d\) from \eqref{eq:defect}, the interval \([a_u,b_u]\)
contains \(\Re f(u)\), and \(j_u\) bounds \(|\Im f(u)|\), where
\[
\begin{aligned}
 A_u&=\bar d(|u\sin u|+u^2),\\
 a_u&=\max\{-1,\cos u-A_u,1-u^2/2\},\\
 b_u&=\min\{1,\cos u+A_u,1-u^2/2+\kappa_+\rho u^3\},\\
 j_u&=\min\{\sqrt{2\bar d}|\sin u-u\cos u|+\bar d u^2,\,
                    u^3\sqrt{\rho^2-\eta^2}/6\}.
\end{aligned}
\]
Combining these component bounds with the symmetrized modulus estimate gives
\(|f(u)|\le F_u\), where
\[
 F_u=\min\{1,\pos{1-2u^2q(2\pi t)}^{1/2},
                    [\max(|a_u|,|b_u|)^2+j_u^2]^{1/2}\}.
\]
The same rectangle bounds the distance to \(\cos u\):
\(|f(u)-\cos u|\le E_u\), with
\[
 e_u=\min\{A_u,\max(|a_u-\cos u|,|b_u-\cos u|)\},\qquad
 E_u=(e_u^2+j_u^2)^{1/2}.
\]

The Gaussian comparison, the triangle inequality and the two-point
comparison now give three bounds on \(|f(u)^n-N(t)|\). Taking their
minimum, define \(M(t)=F_u^n\) and
\[
\begin{aligned}
 R(t)=\min\bigl\{&
 n\rho u^3D(\rho,r,2\pi t/r)\max(F_u,B)^{n-1},\\
 &F_u^n+N,\quad
 nE_u\max(F_u,|\cos u|)^{n-1}+|\cos(u)^n-N|\bigr\}.
\end{aligned}
\]
These bounds remain valid at \(n=1\) and at \(t=0\).
Therefore \(\Delta_n\le\mathcal U_+(n,\rho,r)\), where
\begin{align}
 \mathcal U_+:={}&
 2\int_0^{t_0}|K(t)|R(t)\dd t
 +2\int_{t_0}^1|K(t)|M(t)\dd t \notag\\
 &+2\int_0^{t_0}|K_c(t)|N(t)\dd t
 +\frac1\pi\int_{t_0}^{\infty}N(t)\frac{\dd t}{t},
 \label{eq:U-def}\tag{4.4}
\end{align}
with \(K_c(t)=K(t)-i/(2\pi t)\). Thus \(\mathcal U_+\) bounds the
smoothing functional using only \(n,\rho,r\).

\subsection{Large-sample reduction}

For \(n\ge100\), define
\[
 Q(t)=\frac{(2\pi t)^2q(2\pi t)}{r^2L^2},\qquad
 Q_{\rm rate}(t)=\frac{39(2\pi t)^2}{100r^2L^2},\qquad
 \alpha(L)=\max\{99/100,1-L^2\}.
\]
The two inequalities \(n(99/100)\le n-1\) and
\(n(1-L^2)=n-\rho^2\le n-1\) prove
\[
 n\alpha(L)\le n-1.
 \label{eq:alpha}\tag{4.5}
\]
Since the one-factor envelopes are nonnegative, an estimate
\(\max(F_u,B)\le e^{-q_0}\), \(q_0\ge0\), implies
\[
 \max(F_u,B)^{n-1}\le e^{-\alpha(L)nq_0}.
\]
On \(r\ge19/10\) and \(t\le1/4\), one has
\(u\le5/6\), so Lemma~\ref{lem:radial-rate} applies. Combining it
with Lemma~\ref{lem:q}, put
\[
 Q_s(t)=
 \begin{cases}
 \max\{Q(t),Q_{\rm rate}(t)\},&t\le1/4,\\
 Q(t),&t>1/4.
 \end{cases}
\]
The normalized low-frequency difference integrand is then at most
\[
 \frac{2|K(t)|(2\pi t)^3}{r^3L^3}
 D_*(r,2\pi t/r)e^{-\alpha(L)Q_s(t)}.
 \label{eq:large-new}\tag{4.6}
\]
The estimate \(u\le5/6\) uses only \(\rho\ge1\), \(r\ge19/10\)
and the certified upper enclosure of \(\pi\).

For \(1/16<L\le56/45\), the certificate uses the weaker fixed
coefficient \(99/100\) in \eqref{eq:large-new}. It takes the minimum
of this bound and the Gaussian and triangle bounds obtained using
\(D_0\) and \eqref{eq:modulus}:
\[
 \frac{2|K(t)|}{L}\min\left\{
 \frac{(2\pi t)^3}{r^3L^2}D_0(r,2\pi t/r)e^{-99Q(t)/100},
 \ e^{-Q(t)}+N(t)\right\}.
\]
The high-frequency and correction integrands are bounded, respectively, by
\[
 \frac{2|K(t)|}{L}e^{-Q(t)},\qquad
 \frac{2|K_c(t)|}{L}N(t).
\]
A parameter interval not contained in \(r\ge19/10\) uses the bound with
\(D_0\) throughout.

For \(0<L\le1/16\), the variable coefficient in
\eqref{eq:large-new} is used when \(r\ge19/10\). The substitutions
\(t=Ly\) near zero and \(t=1-Ly\) near one remove the apparent
\(L^{-1}\) singularities before interval evaluation.
Section~\ref{sec:certificate} specifies the endpoint truncation and
the exact allowance for the omitted intervals.

Let \(\mathcal P_n\) denote the original right side of
\eqref{eq:prawitz}, with \(h\) the characteristic function of \(S_n\)
and parameters \eqref{eq:parameters}.
Then \(\Delta_n\le\mathcal P_n\le\mathcal U_+\).
The bounds above reduce the theorem to the following scalar inequality.

\begin{lemma}[Exact continuous-domain bound]\label{lem:numerical}
For every admissible law and integer \(n\ge1\) such that
\[
 1\le\rho\le\frac{56}{45}\sqrt n,\qquad
 1\le r\le1+\frac1\rho,
 \label{eq:num-domain}\tag{4.7}
\]
one has
\[
 \frac{\sqrt n}{\rho}\mathcal P_n<\frac{879}{2000}.
 \label{eq:num-bound}\tag{4.8}
\]
\end{lemma}

The finite certificates use \eqref{eq:U-def}; the large-sample
certificates use the upper integrands just derived. The certificate
is over the entire continuous parameter domain, not a sample of
admissible laws.

\begin{proof}[Proof of Theorem~\ref{thm:main}]
The standardized sum has mean zero and variance one by independence.
If \(L\ge1090/879\), Lemma~\ref{lem:universal} gives
\[
 \Delta_n\le109/200\le(879/2000)L.
\]
Otherwise \(L<1090/879<56/45\). The moment constraints place the law
in \eqref{eq:num-domain}. Lemmas~\ref{lem:prawitz}
and~\ref{lem:numerical} give
\(\Delta_n\le\mathcal P_n<(879/2000)L\), which implies the stated
weak bound.
\end{proof}

\section{Exact verification of the scalar bound}\label{sec:certificate}

The preceding sections reduce the proof to Lemma~\ref{lem:numerical}.  Its
certificate proves an inequality over a continuum of parameters: finite data
describe a covering, and a proved-sound interval checker verifies the desired
bound on every member of that covering.

\subsection{Exact interval bounds}

A precision-48 dyadic interval \([l,h]\) represents
\([l/2^{48},h/2^{48}]\), with integer \(l,h\).
Rational conversion uses floor and ceiling; products use the outward
four-corner hull. The verified operations also include division under
a denominator sign condition, nonnegative square root and integer
powers. Each operation has a containment theorem over the entire
input interval.

For \(e^{-x}\), \(x\ge0\), the checker halves the interval until its
upper endpoint is at most \(1/8\), uses degree-13 and degree-12
alternating Taylor bounds, and reverses the reduction by outward
squaring. The Gaussian correction has a separate degree-25/24
enclosure. Sine and cosine on a full cell are enclosed using
degree-33/32 midpoint polynomials, Lagrange remainder bounds, and
the global Lipschitz bounds for their variation from the midpoint.
Certified kernel bounds use the cotangent-gap expansion.
The constants \(\pi,\kappa,\vartheta\) enter through proved rational
enclosures, not machine floating-point values.

For a partition \(\mathcal P\) of \(J\) into closed intervals, let
\([G](I)\) be an interval containing \(G(t)\) for every \(t\in I\).
Then
\[
 \int_J G(t)\dd t
 \le\sum_{I\in\mathcal P}|I|\,\sup[G](I).
\]
Lean proves this implication with the actual outward-rounded
multiplication and summation order of the checker. The normalized
upper endpoint must be strictly below \(879/2000\), and every
interval-order, sign and analytic side condition must pass. Bisection
uses two closed halves whose union is the parent, so shared endpoints
cannot be omitted.

The improper normal tail is also enclosed. If
\(E_1(x)=\int_x^\infty e^{-v}\dd v/v\), then
\[
 \frac1\pi\int_{t_0}^\infty e^{-T^2t^2/2}\frac{\dd t}{t}
 =\frac{E_1(T^2t_0^2/2)}{2\pi}.
\]
Using \(E_1(x)=e^{-x}\int_0^\infty e^{-s}\dd s/(x+s)\),
the checker applies the convex composite trapezoid bound on
\([0,24]\) with 768 intervals, then bounds the remaining integral by
\(e^{-24}/(x+24)\). Monotonicity replaces a positive interval
parameter \(x\) by its lower endpoint in the upper direction.

\subsection{Exhaustive regimes and endpoint limits}

For each \(1\le n<100\), use
\[
 \rho=1+y\left(\frac{56}{45}\sqrt n-1\right),\qquad
 r=1+\frac z\rho,\qquad (y,z)\in[0,1]^2.
\]
Finite prefix codes describe complete binary subdivisions of this square.
Each accepted leaf supplies a strict interval upper bound on its whole
parameter cell. The union of the closed child cells covers each parent;
the certificate is accepted only when this coverage and every leaf bound hold.

For \(n\ge100\), three closed parameter maps cover the regimes
\[
 0<L\le\frac1{16},\qquad
 \frac1{16}<L\le\frac1{10},\qquad
 \frac1{10}<L\le\frac{56}{45}.
\]
In the first two maps the second coordinate is \(r-1\in[0,1]\).
In the upper map it is \(10\eta/\sqrt n=10L(r-1)\in[0,1]\),
and \(r=1+z/(10L)\). The first coordinate affinely rescales \(L\).
The parameter interval is intersected with the proved constraint
\(r\le2\); this intersection preserves every admissible point.

For the small region, the endpoint substitutions retain \(0\le y\le4\).
The low, middle, high-endpoint and Gaussian-tail omissions have
normalized upper bounds
\[
 \frac{13}{10^{12}},\quad \frac1{10^{12}},\quad
 \frac1{10^{12}},\quad \frac1{10^{12}},
\]
respectively. Their sum is below the allowance \(3/10^{11}\).
These are analytic bounds, proved using
\(e^{-32}\le1/(7\cdot10^{13})\), not unverified truncated tails.
After the substitutions the endpoint expressions have certified
extensions at \(L=0\). This covers arbitrarily small positive \(L\)
without imposing an upper limit on \(n\).

Two small-region trees split \(r-1\) at the exact dyadic
\(253327479039591/2^{48}\), which lies just above \(9/10\).
Their closed intervals meet at that point; equality is assigned to
the lower tree. The equalities \(L=1/16\), \(L=1/10\) and \(L=56/45\)
are included in the corresponding closed certificates.
Thus the finite and large-sample cases exhaust \eqref{eq:num-domain}.

\subsection{Formal verification and reproducibility}

The moment identities, characteristic-function bounds, smoothing theorem,
interval containment and continuous coverage are formalized in Lean 4
\cite{demoura2021lean4} with Mathlib \cite{mathlib2020}.
The accompanying \href{https://github.com/haonan-xiao/iid-berry-esseen}{source
repository} contains the complete theorem, exact certificate data and
reproduction instructions. Its formalization guide maps the mathematical
arguments to the Lean development.

The analytic implications and checker soundness are checked by the
Lean kernel. Concrete Boolean evaluation uses \texttt{native\_decide},
which additionally trusts Lean's native compiler. There are no unproved
placeholders or user-declared axioms. Only the finite certificate data,
not the program that searches for them, enter the proof.

This verifies Lemma~\ref{lem:numerical} on the full continuous domain.
The mathematical improvement comes from retaining the first absolute
moment in the characteristic-function comparisons and sharpening their
use in the smoothing inequality. The certificates
establish that these inequalities attain the stated universal constant.

\section*{AI disclosure}

The proof, Lean formalization, exact numerical certificates, and initial
manuscript were developed autonomously using OpenAI Codex, primarily with
GPT-5.6 Sol. The authors supplied the research question and candidate
materials and coordinated the review and preparation of the note.

\begingroup
\footnotesize
\raggedright
\bibliographystyle{plain}
\bibliography{references}
\endgroup

\end{document}
